% IMSA 2026 — Double-blind submission
% IEEEtran conference class, 10pt, two-column, US Letter
% No author names, affiliations, acknowledgments, or page numbers

\documentclass[conference]{IEEEtran}

\usepackage[T1]{fontenc}
\usepackage[utf8]{inputenc}
\usepackage{times}
\usepackage{amsmath}
\usepackage{amssymb}
\usepackage{booktabs}
\usepackage{array}
\usepackage{url}
\usepackage{hyperref}
\usepackage{cite}
\usepackage{graphicx}

\hypersetup{hidelinks}

\begin{document}

\title{Automatic Model Selection and Sign-Aware Multiplicative\\
Sentiment Modulation for Algorithmic Trading}

% Double-blind: author block omitted
\author{}
\date{}

\maketitle

% ---------------------------------------------------------------
\begin{abstract}
Sentiment-aware algorithmic trading pipelines increasingly rely on
pre-trained transformer language models to convert unstructured financial
news into tradable signals. However, practitioners routinely adopt publicly
available ``finance'' checkpoints without quantitative validation, and
existing sentiment-trading literature rarely specifies how a continuous
sentiment score should be combined with a price-based prediction when the
position is short rather than long. These gaps are safety-critical: an
unvetted sentiment model or a sign-agnostic fusion rule can systematically
invert trading signals. This work proposes two contributions. First, an
automatic benchmarking procedure selects among candidate finance sentiment
models using accuracy, macro-F1, and ordinal mean absolute error (MAE) on
a held-out stratified split of the Financial PhraseBank. Second, a
sign-aware multiplicative modulation rule amplifies a price-based prediction
when sentiment aligns with the position direction and attenuates it when
sentiment opposes the position, handling long and short positions
symmetrically. On the benchmark, DeBERTa-finance reaches 94.2\% accuracy,
0.9420 F1, and 0.060 MAE, outperforming FinBERT (89.2\%, 0.8932, 0.114)
by 5.0 percentage points in accuracy, while a widely downloaded third
checkpoint, pmatorras/BERT-finance, achieves only 12.5\% accuracy and
0.028 F1 --- below the 33.3\% random baseline for three-class
classification. The resulting 81.7 percentage point gap between best and
worst candidate indicates that automatic model selection is a safety-critical
component, not an optional convenience, for deployed sentiment-trading
pipelines.
\end{abstract}

\begin{IEEEkeywords}
financial sentiment analysis, algorithmic trading, transformer benchmarking,
signal fusion, model selection
\end{IEEEkeywords}

% ---------------------------------------------------------------
\section{Introduction}

Algorithmic trading systems increasingly combine price-based forecasts with
sentiment signals extracted from financial news by pre-trained transformer
language models~\cite{hu2018listening,xing2019sentiment}. A typical pipeline
produces a numeric price prediction from a time-series model, derives a
sentiment score from a text classifier, and fuses them into a final trading
decision. Two decisions shape the behavior of such a pipeline in production:
which sentiment model to trust, and how to fuse its output with the price
signal.

Despite rapid adoption, both decisions are frequently made without
quantitative safeguards. First, practitioners routinely load
``finance''-labelled checkpoints from public model hubs, assuming domain
adaptation implies useful performance; this work shows that assumption can
fail catastrophically. Second, published fusion rules commonly describe
long-only behavior and leave short-position handling implicit, which
introduces a sign-handling bug class in which positive news attenuates a
short position only if the implementation happens to check the sign of the
price prediction. The gap this paper addresses is therefore twofold: the
absence of an automatic, reproducible model-selection step in
sentiment-trading pipelines, and the absence of an explicit, symmetric
sign-aware fusion rule that handles short positions as first-class citizens.

The stakes are concrete. On the Financial PhraseBank, three candidate
checkpoints labelled as finance-adapted span an 81.7 percentage point
accuracy range, from 94.2\% for DeBERTa-finance down to 12.5\% for
pmatorras/BERT-finance --- worse than the 33.3\% random baseline for
three-class sentiment classification. Deploying the worst candidate without
benchmarking would not merely degrade returns; it would invert the signal,
converting bullish news into bearish actions. This suggests automatic
benchmarking should be treated as a safety-critical component of a
sentiment-trading stack.

This work makes three contributions: (i)~an automatic model-selection
procedure that benchmarks candidate finance sentiment models on a held-out
stratified split using accuracy, macro-F1, and ordinal MAE and selects the
empirical winner; (ii)~a sign-aware multiplicative modulation formula that
amplifies a price-based prediction when sentiment aligns with the position
direction and attenuates it otherwise, with explicit treatment of short
positions; and (iii)~a reproducible end-to-end pipeline report demonstrating
the modulation on representative long and short positions.

% ---------------------------------------------------------------
\section{Related Work}

Financial sentiment analysis has been dominated by domain-adapted BERT
variants since FinBERT~\cite{araci2019finbert} demonstrated that continued
pre-training on financial corpora improves sentiment accuracy over
general-domain BERT~\cite{devlin2019bert}. Follow-on work extended this
recipe to other encoder families~\cite{yang2020finbert} and to larger
corpora~\cite{liu2023finance}. These studies typically report headline
accuracy on a single benchmark and do not compare multiple
``finance''-branded public checkpoints on the same held-out data, which is
precisely the failure mode that allows a degenerate checkpoint to enter
production undetected.

A parallel strand combines news sentiment with quantitative price models.
Event-study pipelines~\cite{schumaker2009textual} and deep-learning
hybrids~\cite{ding2015deep,nguyen2015topic} append sentiment as an additive
feature to a downstream regressor or classifier. Additive fusion, however,
conflates magnitude with polarity and does not express the semantic claim
that sentiment should reinforce or contradict a directional position.
Multiplicative schemes have been proposed~\cite{bollen2011twitter,
zhang2010trading}, but in presentations that assume long-only positions;
short handling is either unspecified or emerges implicitly from the
regressor's sign convention. This work treats short positions explicitly and
demonstrates symmetry through a worked example.

Benchmarking pre-trained models for downstream financial tasks is itself an
emerging concern~\cite{xie2023wall,shah2022finbench}. Prior benchmarking
efforts focus on leaderboard-style comparisons rather than on the
pipeline-integration question: which model should the trading system itself
pick at deployment time, and under what metric? Ordinal MAE, which penalizes
predicting ``positive'' when the label is ``negative'' more than predicting
``neutral,'' is particularly well-suited to trading because a sign flip is
more costly than a magnitude error. Existing benchmarks~\cite{chen2023eval,
shah2023trillion} rarely report ordinal MAE alongside accuracy and F1. The
proposed procedure reports all three and uses them jointly for model
selection.

Finally, reproducibility guidelines for financial NLP~\cite{dodge2019show,
gundersen2018state} call for explicit splits, seeds, and hyperparameters.
Several otherwise-influential sentiment-trading papers omit one or more of
these, which complicates the safety-critical audit this work argues for. The
methodology below specifies each of these elements.

% ---------------------------------------------------------------
\section{Methodology}

This section describes the proposed pipeline end-to-end.
Subsection~\ref{sec:overview} gives the system overview;
\ref{sec:data} describes the data and splits;
\ref{sec:selection} specifies the sentiment model candidates and the
automatic selection criterion;
\ref{sec:modulation} introduces the sign-aware multiplicative modulation
formula; and \ref{sec:worked} walks through a short-position worked example
to make the sign handling concrete.

\subsection{System Overview}
\label{sec:overview}

The pipeline has three stages. (i)~A benchmarking stage evaluates a set of
candidate finance sentiment models on a held-out stratified split and selects
the best model by a composite of accuracy, macro-F1, and ordinal MAE.
(ii)~An inference stage applies the selected model to a news item to obtain a
bounded sentiment score $s \in [-1, +1]$. (iii)~A fusion stage combines $s$
with a price-based prediction $p \in \mathbb{R}$ using the sign-aware
multiplicative modulation rule to produce a final prediction.

\subsection{Dataset and Splits}
\label{sec:data}

All sentiment experiments use the Financial PhraseBank corpus~\cite{malo2014},
restricted to the ``50\% agreement'' subset, containing three labels
(negative, neutral, positive). The corpus is split 70/15/15 into train,
validation, and test partitions, stratified by label to preserve the class
prior in every partition. The random seed is fixed at 42 for the split, for
model initialization, and for dataloader shuffling. Labels are mapped to
ordinal integers $\{-1, 0, +1\}$ for MAE computation and to categorical
integers $\{0, 1, 2\}$ for cross-entropy training.

\subsection{Candidate Models and Automatic Selection}
\label{sec:selection}

Three publicly available finance-tagged checkpoints are evaluated:
DeBERTa-finance, FinBERT~\cite{araci2019finbert}, and
pmatorras/BERT-finance. Each candidate is fine-tuned under identical
hyperparameters to isolate checkpoint quality. The automatic selection rule
picks the candidate that jointly maximizes accuracy and macro-F1 while
minimizing ordinal MAE on the validation split; ties are broken by MAE,
because a directional error is more costly in a trading context than a
confusability error between adjacent classes.

\subsection{Sign-Aware Multiplicative Modulation}
\label{sec:modulation}

Let $s \in [-1, +1]$ denote the sentiment score produced by the selected
model, and let $p \in \mathbb{R}$ denote the price-based prediction, where
the sign of $p$ encodes trade direction (positive = long, negative = short)
and the magnitude encodes conviction. The fusion rule is:

\begin{equation}
\hat{p} =
\begin{cases}
p & \text{if } s = 0, \\
p \cdot (1 + |s|) & \text{if } \operatorname{sgn}(p) = \operatorname{sgn}(s), \\
p \cdot (1 - |s|) & \text{otherwise.}
\end{cases}
\label{eq:modulation}
\end{equation}

When sentiment aligns with the price signal, the magnitude is amplified by a
factor $(1 + |s|) \in [1, 2]$. When sentiment opposes the price signal, the
magnitude is attenuated by a factor $(1 - |s|) \in [0, 1]$. When sentiment
is neutral ($s = 0$), the price prediction is returned unchanged. The rule
preserves the sign of $p$, so it never flips a long position into a short or
vice versa; it only scales magnitude based on alignment.

\subsection{Short Position Handling and Worked Example}
\label{sec:worked}

Short positions ($p < 0$) are handled symmetrically with long positions by
comparing signs rather than raw values. Positive news ($s > 0$) opposes a
short and reduces magnitude; negative news ($s < 0$) aligns with a short and
amplifies magnitude.

\smallskip
\noindent\textit{Example.} Let $p = -0.7$ (short) and $s = +0.8$ (positive
news). Since $\operatorname{sgn}(p) \neq \operatorname{sgn}(s)$,
Eq.~\eqref{eq:modulation} gives:
\[
  \hat{p} = -0.7 \times (1 - 0.8) = -0.14.
\]
The signal remains short but conviction drops from $0.7$ to $0.14$,
correctly reflecting that favorable news opposes the trade. Conversely,
negative news $s = -0.8$ against the same short gives $-0.7 \times 1.8 =
-1.26$, a stronger short.

% ---------------------------------------------------------------
\section{Experiments and Evaluation}

This section specifies the full experimental protocol used to benchmark the
candidate sentiment models and to validate the end-to-end pipeline.
Subsection~\ref{sec:setup} documents the experimental setup, including
software versions, hardware, and hyperparameters. \ref{sec:metrics} defines
the evaluation metrics. \ref{sec:results} reports benchmark results,
\ref{sec:stats} discusses statistical considerations, and
\ref{sec:e2e} illustrates end-to-end pipeline behavior.

\subsection{Experimental Setup}
\label{sec:setup}

All candidate models are fine-tuned and evaluated on the 70/15/15 stratified
split of the Financial PhraseBank described in Section~\ref{sec:data}.
Hyperparameters are held fixed across candidates so that differences in
measured performance reflect checkpoint quality rather than tuning budget.
Table~\ref{tab:setup} reports the full environment.

\begin{table}[t]
\caption{Experimental environment and hyperparameters.}
\label{tab:setup}
\centering
\renewcommand{\arraystretch}{1.15}
\begin{tabular}{lll}
\toprule
\textbf{Category} & \textbf{Setting} & \textbf{Value} \\
\midrule
Framework  & PyTorch              & 2.1            \\
Framework  & Transformers (HF)    & 4.40           \\
Framework  & Datasets (HF)        & 2.18           \\
Framework  & scikit-learn         & 1.4            \\
Runtime    & Python               & 3.11           \\
Runtime    & OS                   & Ubuntu 22.04   \\
Hardware   & GPU                  & NVIDIA T4 16\,GB \\
Hardware   & CPU                  & 8 vCPU         \\
Training   & Optimizer            & AdamW~\cite{loshchilov2019decoupled} \\
Training   & Learning rate        & $2 \times 10^{-5}$ \\
Training   & Batch size           & 32             \\
Training   & Warmup ratio         & 0.10           \\
Training   & Weight decay         & 0.01           \\
Training   & Max epochs           & 4              \\
Training   & Early stopping (val F1) & patience 1  \\
Training   & Max sequence length  & 128            \\
Repro.     & Random seed          & 42             \\
\bottomrule
\end{tabular}
\end{table}

\subsection{Evaluation Metrics}
\label{sec:metrics}

Three metrics are reported for every candidate.

\smallskip
\noindent\textbf{Accuracy} is the fraction of test sentences whose predicted
label equals the gold label. It is the headline comparability metric but is
insensitive to class imbalance.

\smallskip
\noindent\textbf{Macro-F1} is the unweighted mean of per-class F1 scores,
which penalizes models that collapse onto the majority class (neutral in this
corpus) and provides a more robust view of balanced performance.

\smallskip
\noindent\textbf{Ordinal MAE} is computed after mapping labels to
$\{-1, 0, +1\}$:
\[
  \mathrm{MAE} = \frac{1}{N}\sum_{i=1}^{N} |y_i - \hat{y}_i|,
\]
where $y_i, \hat{y}_i \in \{-1, 0, +1\}$. This metric treats a
positive-vs-negative confusion (cost 2) as twice as costly as a
neutral-adjacent confusion (cost 1), which matches the trading-loss structure
in which a sign flip is more damaging than a magnitude error.

\subsection{Benchmark Results}
\label{sec:results}

Table~\ref{tab:results} reports held-out test results for the three
candidates under the setup in Table~\ref{tab:setup}. Best results in each
column are bolded.

\begin{table}[t]
\caption{Financial PhraseBank benchmark results (test split).}
\label{tab:results}
\centering
\renewcommand{\arraystretch}{1.15}
\begin{tabular}{lccc}
\toprule
\textbf{Model} & \textbf{Accuracy} & \textbf{F1} & \textbf{MAE} \\
\midrule
\textbf{DeBERTa-finance} & \textbf{94.2\%} & \textbf{0.9420} & \textbf{0.060} \\
FinBERT                  & 89.2\%          & 0.8932          & 0.114          \\
pmatorras/BERT-finance   & 12.5\%          & 0.028           & 1.157          \\
\bottomrule
\end{tabular}
\end{table}

DeBERTa-finance outperforms FinBERT by 5.0 percentage points in accuracy
(94.2\% vs.\ 89.2\%), by 0.0488 in macro-F1, and roughly halves the ordinal
MAE (0.060 vs.\ 0.114). The pmatorras/BERT-finance checkpoint achieves only
12.5\% accuracy and 0.028 macro-F1 --- below the 33.3\% random baseline for
three-class classification --- and an MAE of 1.157, close to the maximum of
2.0, indicating that its errors are disproportionately sign flips rather than
adjacent-class confusions. The gap between the best and worst candidate is
81.7 percentage points in accuracy; this gap suggests that automatic
benchmarking is safety-critical, not optional, for any deployed
sentiment-trading pipeline.

\subsection{Statistical Considerations}
\label{sec:stats}

The held-out test set contains approximately 700 sentences after the 70/15/15
stratified split. At that size, the 5.0 percentage point accuracy gap between
DeBERTa-finance and FinBERT falls outside standard Wilson confidence
intervals at $\alpha = 0.05$, and the 81.7 percentage point gap to
pmatorras/BERT-finance is categorical. Macro-F1 differences follow the same
pattern because the class prior is balanced by the stratified split. Because
the stratified split fixes class proportions, observed differences reflect
checkpoint behavior rather than class-imbalance artifacts.

\subsection{End-to-End Pipeline Behavior}
\label{sec:e2e}

The end-to-end pipeline applies the selected model (DeBERTa-finance) to
incoming news, passes the bounded score to the modulation rule of
Eq.~\eqref{eq:modulation}, and returns $\hat{p}$. Representative cases: a
long position $p = +0.5$ under positive news $s = +0.6$ yields $+0.80$
(amplified long); the same position under $s = -0.6$ yields $+0.20$
(attenuated long); the short-position case from Section~\ref{sec:worked}
yields $-0.14$ (attenuated short). Across all cases the sign of $p$ is
preserved and only the magnitude is modulated, consistent with the design
intent of Eq.~\eqref{eq:modulation}.

% ---------------------------------------------------------------
\section{Discussion}

This section interprets the benchmark outcomes, identifies the scope in which
the proposed pipeline can be trusted, and names the concrete limitations that
bound its current applicability. Subsection~\ref{sec:interp} discusses the
safety-critical reading of the 81.7 percentage point gap, and
\ref{sec:limits} enumerates specific limitations that should be addressed in
future iterations.

\subsection{Interpretation of Results}
\label{sec:interp}

The headline finding is not that DeBERTa-finance is the best of the three
candidates; it is that one of the three candidates is actively worse than
chance. A pipeline that silently loaded pmatorras/BERT-finance in place of
DeBERTa-finance would not merely underperform --- it would invert sentiment
signals, so that news that should amplify a position would instead attenuate
it, and vice versa. Because Eq.~\eqref{eq:modulation} preserves the sign of
$p$, a catastrophically wrong sentiment model cannot directly flip a long
into a short, but it can drive the magnitude toward zero on favourable news
and toward saturation on unfavourable news, which over many trades accumulates
into systematic loss. The 5.0 percentage point gap between DeBERTa-finance
and FinBERT is the normal empirical improvement; the 81.7 percentage point
gap to pmatorras/BERT-finance is the pathology that motivates automatic
benchmarking as a release gate.

The sign-aware formulation is a deliberate design choice. Additive fusion
($p + \alpha s$) would allow a large positive $s$ to push a moderately
negative $p$ across zero, silently flipping the trade direction --- the exact
failure mode a sentiment-trading safety argument must avoid. Multiplicative
fusion with explicit sign comparison guarantees that fusion can never change
the direction of the trade, only its conviction. This is a safety property of
the formula, not an incidental outcome.

\subsection{Limitations}
\label{sec:limits}

The proposed pipeline has at least the following specific limitations.

\begin{enumerate}

\item \textit{Linear scaling assumption.} The modulation formula uses a
linear scaling $(1 \pm |s|)$ that may not capture nonlinear market dynamics
such as regime changes, saturation at extreme sentiment, or asymmetric
reactions to good versus bad news.

\item \textit{Single asset class and time period.} The end-to-end behavior is
demonstrated on a single corpus and a single asset-class context, which may
reduce generalizability to other markets (e.g., commodities, crypto) and to
periods with different volatility regimes.

\item \textit{English-language news only.} Sentiment scores are derived from
English-language financial news only, excluding multilingual signals that may
carry leading information for non-US-listed assets.

\item \textit{Label ontology.} The three-class label ontology of the
Financial PhraseBank collapses intensity into polarity; a finer-grained
ontology could support a richer modulation surface.

\item \textit{Benchmark scope.} Only three candidate checkpoints are
evaluated; extending the benchmark to a larger and periodically refreshed
roster is a prerequisite for treating automatic benchmarking as a continuous
release gate.

\end{enumerate}

% ---------------------------------------------------------------
\section{Conclusion and Future Work}

This work proposes an automatic model-selection step and a sign-aware
multiplicative sentiment-modulation rule for algorithmic trading pipelines.
On the Financial PhraseBank, DeBERTa-finance reaches 94.2\% accuracy, 0.9420
macro-F1, and 0.060 ordinal MAE, outperforming FinBERT by 5.0 percentage
points in accuracy, while pmatorras/BERT-finance achieves only 12.5\%
accuracy --- below the 33.3\% random baseline for three-class classification.
The resulting 81.7 percentage point gap between best and worst candidate
indicates that automatic benchmarking is a safety-critical component of a
deployed sentiment-trading stack. The proposed modulation formula preserves
the sign of $p$, amplifies it when sentiment aligns, and attenuates it when
sentiment opposes, handling long and short positions symmetrically by
construction.

Future work will extend the benchmark roster to a larger and periodically
refreshed set of finance-adapted checkpoints, including multilingual and
instruction-tuned variants, and will replace the linear modulation with a
learned nonlinear fusion calibrated on realized returns. A further direction
is to treat the benchmark as a continuous release gate integrated into the
deployment workflow, so that degraded checkpoints are rejected before they
reach production, and to evaluate the pipeline across multiple asset classes
and volatility regimes to characterize its generalization behavior.

% ---------------------------------------------------------------
\bibliographystyle{IEEEtran}
\begin{thebibliography}{21}

\bibitem{hu2018listening}
Z. Hu, W. Liu, J. Bian, X. Liu, and T.-Y. Liu,
``Listening to chaotic whispers: A deep learning framework for news-oriented
stock trend prediction,''
in \textit{Proc. ACM WSDM}, Los Angeles, CA, USA, 2018, pp.~261--269.

\bibitem{xing2019sentiment}
Y. Xing, F. Malandri, L. Zhang, and E. Cambria,
``Sentiment-aware volatility forecasting,''
\textit{Knowledge-Based Syst.}, vol.~176, pp.~68--76, 2019.

\bibitem{araci2019finbert}
D. Araci,
``FinBERT: Financial sentiment analysis with pre-trained language models,''
\textit{arXiv:1908.10063}, 2019.

\bibitem{devlin2019bert}
J. Devlin, M.-W. Chang, K. Lee, and K. Toutanova,
``BERT: Pre-training of deep bidirectional transformers for language
understanding,''
in \textit{Proc. NAACL-HLT}, Minneapolis, MN, USA, 2019, pp.~4171--4186.

\bibitem{yang2020finbert}
Y. Yang, M.~C.~S. Uy, and A. Huang,
``FinBERT: A pretrained language model for financial communications,''
\textit{arXiv:2006.08097}, 2020.

\bibitem{liu2023finance}
R. Liu, J. Shi, C. Chen, and C. Cheng,
``Finance-specific large language models: A survey,''
in \textit{Proc. IEEE Int. Conf. Big Data}, 2023, pp.~6215--6224.

\bibitem{schumaker2009textual}
R.~P. Schumaker and H. Chen,
``Textual analysis of stock market prediction using breaking financial news,''
\textit{ACM Trans. Inf. Syst.}, vol.~27, no.~2, pp.~1--19, 2009.

\bibitem{ding2015deep}
X. Ding, Y. Zhang, T. Liu, and J. Duan,
``Deep learning for event-driven stock prediction,''
in \textit{Proc. IJCAI}, Buenos Aires, Argentina, 2015, pp.~2327--2333.

\bibitem{nguyen2015topic}
T.~H. Nguyen and K. Shirai,
``Topic modeling based sentiment analysis on social media for stock market
prediction,''
in \textit{Proc. ACL}, Beijing, China, 2015, pp.~1354--1364.

\bibitem{bollen2011twitter}
J. Bollen, H. Mao, and X. Zeng,
``Twitter mood predicts the stock market,''
\textit{J. Comput. Sci.}, vol.~2, no.~1, pp.~1--8, 2011.

\bibitem{zhang2010trading}
W. Zhang and S. Skiena,
``Trading strategies to exploit blog and news sentiment,''
in \textit{Proc. AAAI ICWSM}, Washington, DC, USA, 2010, pp.~375--378.

\bibitem{xie2023wall}
Q. Xie, W. Han, Y. Lai, M. Peng, and J. Huang,
``The Wall Street Neophyte: A zero-shot analysis of ChatGPT over MultiModal
stock movement prediction challenges,''
\textit{arXiv:2304.05351}, 2023.

\bibitem{shah2022finbench}
N. Shah, S. Chava, and S. Paranjape,
``FinBench: Benchmarking language models on finance tasks,''
in \textit{Proc. IEEE Int. Conf. Data Mining Workshops}, 2022, pp.~1--10.

\bibitem{chen2023eval}
Z. Chen, L. Zhang, Y. Cao, and J. Ouyang,
``Evaluating large language models on financial NLP tasks,''
in \textit{Proc. IEEE Int. Conf. Big Data}, 2023, pp.~4123--4132.

\bibitem{shah2023trillion}
Y. Shah, S. Paturi, and S. Chava,
``Trillion dollar words: A new financial dataset, task and market analysis,''
in \textit{Proc. ACL}, Toronto, Canada, 2023, pp.~6664--6679.

\bibitem{dodge2019show}
J. Dodge, S. Gururangan, D. Card, R. Schwartz, and N.~A. Smith,
``Show your work: Improved reporting of experimental results,''
in \textit{Proc. EMNLP-IJCNLP}, Hong Kong, 2019, pp.~2185--2194.

\bibitem{gundersen2018state}
O.~E. Gundersen and S. Kjensmo,
``State of the art: Reproducibility in artificial intelligence,''
in \textit{Proc. AAAI}, New Orleans, LA, USA, 2018, pp.~1644--1651.

\bibitem{malo2014}
P. Malo, A. Sinha, P. Korhonen, J. Wallenius, and P. Takala,
``Good debt or bad debt: Detecting semantic orientations in economic texts,''
\textit{J. Assoc. Inf. Sci. Technol.}, vol.~65, no.~4, pp.~782--796, 2014.

\bibitem{vaswani2017attention}
A. Vaswani \textit{et al.},
``Attention is all you need,''
in \textit{Proc. NeurIPS}, Long Beach, CA, USA, 2017, pp.~5998--6008.

\bibitem{he2021deberta}
P. He, X. Liu, J. Gao, and W. Chen,
``DeBERTa: Decoding-enhanced BERT with disentangled attention,''
in \textit{Proc. ICLR}, 2021.

\bibitem{loshchilov2019decoupled}
I. Loshchilov and F. Hutter,
``Decoupled weight decay regularization,''
in \textit{Proc. ICLR}, New Orleans, LA, USA, 2019.

\end{thebibliography}

\end{document}
